Large language models (LLMs) are increasingly used as reasoning evaluators, yet whether they can reliably assess their own logical reasoning remains poorly understood. We investigate whether LLMs exhibit asymmetric detection rates for logical fallacies when evaluating reasoning attributed to themselves versus other models. Using 195 fallacious reasoning examples spanning 13 categories from the \logicdataset dataset, we conduct three experiments with \gptfull and \gptmini: (1) an attribution framing experiment presenting identical fallacious text under self-attributed, other-attributed, and neutral conditions; (2) a cross-model evaluation where models assess both their own and each other's generated reasoning; and (3) a prompting intervention testing whether explicit fallacy-checking instructions reduce bias. Contrary to our hypothesis, \gptfull does not exhibit statistically significant self-evaluation bias: the attribution bias is only $+3.1\%$ ($p = 0.146$, McNemar's test, Cohen's $h = 0.129$). In genuine cross-model evaluation, self-evaluation accuracy is slightly \emph{higher} than cross-evaluation ($82.0\%$ vs.\ $79.5\%$). However, we identify category-specific asymmetries up to $13.3\%$ for false dilemma, faulty generalization, and intentional fallacies. Our results suggest that modern frontier LLMs are remarkably robust to attribution framing manipulations, though category-specific blindspots persist and merit investigation at larger scale.
